%% This is file `TBench-template.tex'.
%%
%%
%% This file is part of the 'TBench-template Bundle'.
%% BenchCouncil Transactions on Benchmarks, Standards and Evaluations (TBench)
%%
%%
%% This class file was adapted from a template originally created by Peter Jones for the journal "Microscopy and Microanalysis", available at:
%% https://www.overleaf.com/latex/templates/microscopy-and-microanalysis-template/ndmcfsytwfbm
%% Licensed under the Creative Commons Attribution 4.0 International (CC BY 4.0).
%% [https://creativecommons.org/licenses/by/4.0/]
%% Modifications and adaptations have been made to serve the formatting requirements of TBench. 
%%
%% This derived class is proprietary and intended for use in submissions to TBench only.
%%
%% The list of all files belonging to the 'TBench-template Bundle' is
%% provided in the file `manifest.txt`.
%%
%% Copyright © 2025 BenchCouncil Press on Behalf of International Open Benchmark Council. All rights reserved.
%% 
%%
%% The list of all files belonging to the 'TBench-template Bundle' is
%% given in the file `manifest.txt'.
%%
%% Template article for 'BenchCouncil Transactions on Benchmarks, Standards and Evaluations' document class `TBench-template'
%% with bibliographic references
%%
%% Version 1.0, 2024.

\documentclass[namedate,webpdf,contemporary,large]{TBench-template}%
\usepackage[english]{babel}
\usepackage{booktabs}
\usepackage{adjustbox}
\usepackage{graphicx}
 % Modify the color of serial number indexes
\usepackage{xcolor}
\makeatletter % Use internal commands
% Custom modifications for natbib
\AtBeginDocument{
    \renewcommand*{\citenumfont}{\color{jnlclr}}
}
\makeatother

%\usepackage{showframe}

\graphicspath{{Fig/}}

% line numbers
%\usepackage[mathlines, switch]{lineno}
%\usepackage[right]{lineno}

\theoremstyle{thmstyleone}%
\newtheorem{theorem}{Theorem}%  meant for continuous numbers
%%\newtheorem{theorem}{Theorem}[section]% meant for sectionwise numbers
%% optional argument [theorem] produces theorem numbering sequence instead of independent numbers for Proposition
\newtheorem{proposition}[theorem]{Proposition}%
%%\newtheorem{proposition}{Proposition}% to get separate numbers for theorem and proposition etc.
\theoremstyle{thmstyletwo}%
\newtheorem{example}{Example}%
\newtheorem{remark}{Remark}%
\theoremstyle{thmstylethree}%
\newtheorem{definition}{Definition}

\begin{document}

\journaltitle{BenchCouncil Transactions on Benchmarks, Standards and Evaluations}
\DOI{DOI HERE}
\copyrightyear{2026}
\pubyear{2026}
%\access{Advance Access Publication Date: Day Month Year}
\appnotes{Original Article}

\firstpage{1}

%\subtitle{Subject Section}
\title[Short Article Title]{Article Title}

\author[1,$\ast$]{First Author}
\author[2]{Second Author}
\author[3]{Third Author}
\author[3]{Fourth Author}
\author[4]{Fifth Author\ORCID{0000-0000-0000-0000}}

\authormark{Author Name et al.}

\address[1]{\orgdiv{Department}, \orgname{Organization}, \orgaddress{\street{Street}, \postcode{Postcode}, \state{State}, \country{Country}}}
\address[2]{\orgdiv{Department}, \orgname{Organization}, \orgaddress{\street{Street}, \postcode{Postcode}, \state{State}, \country{Country}}}
\address[3]{\orgdiv{Department}, \orgname{Organization}, \orgaddress{\street{Street}, \postcode{Postcode}, \state{State}, \country{Country}}}
 
\address[4]{\orgdiv{Department}, \orgname{Organization}, \orgaddress{\street{Street}, \postcode{Postcode}, \state{State}, \country{Country}}}

\corresp[$\ast$]{Corresponding author. \href{mailto:email-id.com}{email-id.com}}

\received{Date}{0}{Year}
%\revised{Date}{0}{Year}
\accepted{Date}{0}{Year}

\abstract{Abstracts must be able to stand alone and so cannot contain citations to
the paper's references, equations, etc. An abstract must consist of a single
paragraph and be concise. Because of online formatting, abstracts must appear
as plain as possible.}
\keywords{Keyword1, Keyword2, Keyword3, Keyword4}

% \boxedtext{
% \begin{itemize}
% \item Key boxed text here.
% \item Key boxed text here.
% \item Key boxed text here.
% \end{itemize}}

\maketitle

\section{Introduction}
The introduction introduces the context and summarizes the manuscript. It is importantly to clearly state the contributions of this piece of work. The introduction introduces the context and summarizes the manuscript. It is importantly to clearly state the contributions 
of this piece of work. The introduction introduces the context and summarizes the manuscript. It is importantly to clearly state the contributions 
of this piece of work. 

This is an example of a new parapgraph with a numbered footnote\footnote{\url{https://www.benchcouncil.org/}} and a second footnote marker.\footnote{Example of footnote text.}



\section{This is an example for first level head - section head}\label{sec2}
Dorem falis untera mocelit, vandor ipsum claret nodis prelunix. Sintora velit ad nume dravicula sentor, quamin elit sed ornare pellivus. Curnis montel avora sit amet quistrum denari, vivamus tempor lacus vel urna commodo. Ameton visco relit nulla facelis, quis velit interdum coram purus. Netram felis euismod grana trentor, nulla lobortis sapien est (refer Section~\ref{sec5}).

\subsection{This is an example for second level head - subsection head}\label{subsec1}

Dorem falis untera mocelit, vandor ipsum claret nodis prelunix. Sintora velit ad nume dravicula sentor, quamin elit sed ornare pellivus. Curnis montel avora sit amet quistrum denari, vivamus tempor lacus vel urna commodo. Ameton visco relit nulla facelis, quis velit interdum coram purus. Netram felis euismod grana trentor, nulla lobortis sapien est. 

\subsubsection{This is an example for third level head - subsubsection head}\label{subsubsec1}

Dorem falis untera mocelit, vandor ipsum claret nodis prelunix. Sintora velit ad nume dravicula sentor, quamin elit sed ornare pellivus. Curnis montel avora sit amet quistrum denari, vivamus tempor lacus vel urna commodo. Ameton visco relit nulla facelis, quis velit interdum coram purus. Netram felis euismod grana trentor, nulla lobortis sapien est. Vistula nec tenetur dis placarat, justo varius fermentum nulla eget. Nullam rhoncus doram nec felis convallis, placerat nunc quamus tellor. %Duis aute irure dolor in reprehenderit in voluptate velit esse cillum dolore eu fugiat nulla pariatur. Excepteur sint occaecat cupidatat non proident, sunt in culpa qui officia deserunt mollit anim id est laborum.

\paragraph{This is an example for fourth level head - paragraph head}

Dorem falis untera mocelit, vandor ipsum claret nodis prelunix. Sintora velit ad nume dravicula sentor, quamin elit sed ornare pellivus. Curnis montel avora sit amet quistrum denari, vivamus tempor lacus vel urna commodo. Ameton visco relit nulla facelis, quis velit interdum coram purus. Netram felis euismod grana trentor, nulla lobortis sapien est. Vistula nec tenetur dis placarat, justo varius fermentum nulla eget. Nullam rhoncus doram nec felis convallis, placerat nunc quamus tellor.

\section{This is an example for first level head}\label{sec3}

\subsection{This is an example for second level head - subsection head}\label{subsec2}

\subsubsection{This is an example for third level head - subsubsection head}\label{subsubsec2}

Dorem falis untera mocelit, vandor ipsum claret nodis prelunix. Sintora velit ad nume dravicula sentor, quamin elit sed ornare pellivus. Curnis montel avora sit amet quistrum denari, vivamus tempor lacus vel urna commodo. Ameton visco relit nulla facelis, quis velit interdum coram purus. Netram felis euismod grana trentor, nulla lobortis sapien est. Vistula nec tenetur dis placarat, justo varius fermentum nulla eget. Nullam rhoncus doram nec felis convallis, placerat nunc quamus tellor.

\paragraph{This is an example for fourth level head - paragraph head}

Dorem falis untera mocelit, vandor ipsum claret nodis prelunix. Sintora velit ad nume dravicula sentor, quamin elit sed ornare pellivus. Curnis montel avora sit amet quistrum denari, vivamus tempor lacus vel urna commodo. Ameton visco relit nulla facelis, quis velit interdum coram purus. Netram felis euismod grana trentor, nulla lobortis sapien est. Vistula nec tenetur dis placarat, justo varius fermentum nulla eget. Nullam rhoncus doram nec felis convallis, placerat nunc quamus tellor.


\section{Equations}\label{sec4}

Equations in \LaTeX{} can either be inline or set as display equations. For
inline equations use the \verb+$...$+ commands. Eg: the equation
$H\psi = E \psi$ is written via the command \verb+$H \psi = E \psi$+.

For display equations (with auto generated equation numbers)
one can use the equation or eqnarray environments:
\begin{equation}
\|\tilde{X}(k)\|^2 \leq\frac{\sum\limits_{i=1}^{p}\left\|\tilde{Y}_i(k)\right\|^2+\sum\limits_{j=1}^{q}\left\|\tilde{Z}_j(k)\right\|^2 }{p+q},\label{eq1}
\end{equation}
where,
\begin{align}
D_\mu &=  \partial_\mu - ig \frac{\lambda^a}{2} A^a_\mu \nonumber \\
F^a_{\mu\nu} &= \partial_\mu A^a_\nu - \partial_\nu A^a_\mu + g f^{abc} A^b_\mu A^a_\nu.\label{eq2}
\end{align}
Notice the use of \verb+\nonumber+ in the align environment at the end
of each line, except the last, so as not to produce equation numbers on
lines where no equation numbers are required. The \verb+\label{}+ command
should only be used at the last line of an align environment where
\verb+\nonumber+ is not used.
\begin{equation}
Y_\infty = \left( \frac{m}{\textrm{GeV}} \right)^{-3}
    \left[ 1 + \frac{3 \ln(m/\textrm{GeV})}{15}
    + \frac{\ln(c_2/5)}{15} \right].
\end{equation}
The class file also supports the use of \verb+\mathbb{}+, \verb+\mathscr{}+ and
\verb+\mathcal{}+ commands. As such \verb+\mathbb{R}+, \verb+\mathscr{R}+
and \verb+\mathcal{R}+ produces $\mathbb{R}$, $\mathscr{R}$ and $\mathcal{R}$
respectively (refer Subsubsection~\ref{subsubsec3}).


Dorem falis untera mocelit, vandor ipsum claret nodis prelunix. Sintora velit ad nume dravicula sentor, quamin elit sed ornare pellivus. Curnis montel avora sit amet quistrum denari, vivamus tempor lacus vel urna commodo. Ameton visco relit nulla facelis, quis velit interdum coram purus. Netram felis euismod grana trentor, nulla lobortis sapien est. Vistula nec tenetur dis placarat, justo varius fermentum nulla eget. Nullam rhoncus doram nec felis convallis, placerat nunc quamus tellor.Dorem falis untera mocelit, vandor ipsum claret nodis prelunix. Sintora velit ad nume dravicula sentor, quamin elit sed ornare pellivus. Curnis montel avora sit amet quistrum denari, vivamus tempor lacus vel urna commodo. Ameton visco relit nulla facelis, quis velit interdum coram purus. Netram felis euismod grana trentor, nulla lobortis sapien est. Vistula nec tenetur dis placarat, justo varius fermentum nulla eget. Nullam rhoncus doram nec felis convallis, placerat nunc quamus tellor.


\section{Tables}\label{sec5}

Tables can be inserted via the normal table and tabular environment. To put
footnotes inside tables one has to 
Dorem falis untera mocelit, vandor ipsum claret nodis prelunix. Sintora velit ad nume dravicula sentor, quamin elit sed ornare pellivus. Curnis montel avora sit amet quistrum denari, vivamus tempor lacus vel urna commodo. Ameton visco relit nulla facelis, quis velit interdum coram purus. Netram felis euismod grana trentor, nulla lobortis sapien est. Vistula nec tenetur dis placarat, justo varius fermentum nulla eget. Nullam rhoncus doram nec felis convallis, placerat nunc quamus tellor. 
use the additional ``tablenotes" environment
enclosing the tabular environment. The footnote appears just below the table
itself (refer Tables~\ref{tab1} and \ref{tab2}).


\begin{verbatim}
\begin{table}[t]
\begin{center}
\begin{minipage}{<width>}
\caption{<table-caption>\label{<table-label>}}%
\begin{tabular}{@{}llll@{}}
\toprule
column 1 & column 2 & column 3 & column 4\\
\midrule
row 1 & data 1 & data 2          & data 3 \\
row 2 & data 4 & data 5$^{1}$ & data 6 \\
row 3 & data 7 & data 8      & data 9$^{2}$\\
\botrule
\end{tabular}
\begin{tablenotes}%
\item Source: Example for source.
\item[$^{1}$] Example for a 1st table footnote.
\item[$^{2}$] Example for a 2nd table footnote.
\end{tablenotes}
\end{minipage}
\end{center}
\end{table}
\end{verbatim}


Lengthy tables which do not fit within textwidth should be set as rotated tables. For this, we need to use \verb+\begin{sidewaystable}...+ \verb+\end{sidewaystable}+ instead of\break \verb+\begin{table}...+ \verb+\end{table}+ environment.


\begin{table}[!t]
\caption{Caption text\label{tab1}}%
\begin{tabular*}{\columnwidth}{@{\extracolsep\fill}llll@{\extracolsep\fill}}
\toprule
column 1 & column 2  & column 3 & column 4\\
\midrule
row 1    & data 1   & data 2  & data 3  \\
row 2    & data 4   & data 5$^{1}$  & data 6  \\
row 3    & data 7   & data 8  & data 9$^{2}$  \\
\botrule
\end{tabular*}
\begin{tablenotes}%
\item Source: This is an example of table footnote this is an example of table footnote this is an example of table footnote this is an example of~table footnote this is an example of table footnote
\item[$^{1}$] Example for a first table footnote.
\item[$^{2}$] Example for a second table footnote.
\end{tablenotes}
\end{table}

\begin{table*}[t]
\caption{Example of a lengthy table which is set to full textwidth.\label{tab2}}
\tabcolsep=0pt%%
\begin{tabular*}{\textwidth}{@{\extracolsep{\fill}}lcccccc@{\extracolsep{\fill}}}
\toprule%
& \multicolumn{3}{@{}c@{}}{Element 1$^{1}$} & \multicolumn{3}{@{}c@{}}{Element 2$^{2}$} \\
\cline{2-4}\cline{5-7}%
Project & Energy & $\sigma_{calc}$ & $\sigma_{expt}$ & Energy & $\sigma_{calc}$ & $\sigma_{expt}$ \\
\midrule
Element 3  & 990 A & 1168 & $1547\pm12$ & 780 A & 1166 & $1239\pm100$\\
Element 4  & 500 A & 961  & $922\pm10$  & 900 A & 1268 & $1092\pm40$\\
\botrule
\end{tabular*}
\begin{tablenotes}%
\item Note: This is an example of table footnote this is an example of table footnote this is an example of table footnote this is an example of~table footnote this is an example of table footnote
\item[$^{1}$] Example for a first table footnote.
\item[$^{2}$] Example for a second table footnote.\vspace*{6pt}
\end{tablenotes}
\end{table*}

\begin{sidewaystable}%[!p]
\caption{Tables which are too long to fit, should be written using the ``sidewaystable" environment as shown here\label{tab3}}
\tabcolsep=0pt%
\begin{tabular*}{\textwidth}{@{\extracolsep{\fill}}lcccccc@{\extracolsep{\fill}}}
\toprule%
& \multicolumn{3}{@{}c@{}}{Element 1$^{1}$}& \multicolumn{3}{@{}c@{}}{Element$^{2}$} \\
\cline{2-4}\cline{5-7}%
Projectile & Energy     & $\sigma_{calc}$ & $\sigma_{expt}$ & Energy & $\sigma_{calc}$ & $\sigma_{expt}$ \\
\midrule
Element 3 & 990 A & 1168 & $1547\pm12$ & 780 A & 1166 & $1239\pm100$ \\
Element 4 & 500 A & 961  & $922\pm10$  & 900 A & 1268 & $1092\pm40$ \\
\botrule
\end{tabular*}
\begin{tablenotes}%
\item Note: This is an example of a table footnote this is an example of a table footnote this is an example of a table footnote this is an example of a table footnote this is an example of a table footnote
\item[$^{1}$] This is an example of a table footnote
\end{tablenotes}
\end{sidewaystable}


\section{Figures}\label{sec6}

As per display \LaTeX\ standards one has to use eps images for \verb+latex+ compilation and \verb+pdf/jpg/png+ images for
\verb+pdflatex+ compilation. This is one of the major differences between \verb+latex+
and \verb+pdflatex+. The images should be single-page documents. The command for inserting images
for \verb+latex+ and \verb+pdflatex+ can be generalized. The package used to insert images in \verb+latex/pdflatex+ is the
graphicx package. Figures can be inserted via the normal figure environment as shown in the below example:


\begin{figure}[!t]%
\centering
{\color{black!20}\rule{213pt}{37pt}}
\caption{This is a widefig. This is an example of a long caption this is an example of a long caption  this is an example of a long caption this is an example of a long caption}\label{fig1}
\end{figure}

\begin{figure*}[!t]%
\centering
{\color{black!20}\rule{438pt}{74pt}}
\caption{This is a widefig. This is an example of a long caption this is an example of a long caption  this is an example of a long caption this is an example of a long caption}\label{fig2}
\end{figure*}


\begin{verbatim}
\begin{figure}[t]
        \centering\includegraphics{<eps-file>}
        \caption{<figure-caption>}
        \label{<figure-label>}
\end{figure}
\end{verbatim}

Test text here.

For sample purposes, we have included the width of images in the
optional argument of \verb+\includegraphics+ tag. Please ignore this.
Lengthy figures which do not fit within textwidth should be set in rotated mode. For rotated figures, we need to use \verb+\begin{sidewaysfigure}+ \verb+...+ \verb+\end{sidewaysfigure}+ instead of the \verb+\begin{figure}+ \verb+...+ \verb+\end{figure}+ environment.

\begin{sidewaysfigure}%
\centering
{\color{black!20}\rule{610pt}{102pt}}
\caption{This is an example for a sideways figure. This is an example of a long caption this is an example of a long caption  this is an example of a long caption this is an example of a long caption}\label{fig3}
\end{sidewaysfigure}



\section{Algorithms, Program codes and Listings}\label{sec7}

Packages \verb+algorithm+, \verb+algorithmicx+ and \verb+algpseudocode+ are used for setting algorithms in latex.
For this, one has to use the below format:


\begin{verbatim}
\begin{algorithm}
\caption{<alg-caption>}\label{<alg-label>}
\begin{algorithmic}[1]
. . .
\end{algorithmic}
\end{algorithm}
\end{verbatim}


You may need to refer to the above-listed package documentations for more details before setting an \verb+algorithm+ environment.
To set program codes, one has to use the \verb+program+ package. We need to use the \verb+\begin{program}+ \verb+...+
\verb+\end{program}+ environment to set program codes.

\begin{algorithm}[!t]
\caption{Calculate $y = x^n$}\label{algo1}
\begin{algorithmic}[1]
\Require $n \geq 0 \vee x \neq 0$
\Ensure $y = x^n$
\State $y \Leftarrow 1$
\If{$n < 0$}
        \State $X \Leftarrow 1 / x$
        \State $N \Leftarrow -n$
\Else
        \State $X \Leftarrow x$
        \State $N \Leftarrow n$
\EndIf
\While{$N \neq 0$}
        \If{$N$ is even}
            \State $X \Leftarrow X \times X$
            \State $N \Leftarrow N / 2$
        \Else[$N$ is odd]
            \State $y \Leftarrow y \times X$
            \State $N \Leftarrow N - 1$
        \EndIf
\EndWhile
\end{algorithmic}
\end{algorithm}

Similarly, for \verb+listings+, one has to use the \verb+listings+ package. The \verb+\begin{lstlisting}+ \verb+...+ \verb+\end{lstlisting}+ environment is used to set environments similar to the \verb+verbatim+ environment. Refer to the \verb+lstlisting+ package documentation for more details on this.


\begin{minipage}{\hsize}%
\lstset{language=Pascal}% Set your language (you can change the language for each code-block optionally)
\begin{lstlisting}[frame=single,framexleftmargin=-1pt,framexrightmargin=-17pt,framesep=12pt,linewidth=0.98\textwidth]
for i:=maxint to 0 do
begin
{ do nothing }
end;
Write('Case insensitive ');
Write('Pascal keywords.');
\end{lstlisting}
\end{minipage}


\section{Cross referencing}\label{sec8}

Environments such as figure, table, equation, and align can have a label
declared via the \verb+\label{#label}+ command. For figures and table
environments one should use the \verb+\label{}+ command inside or just
below the \verb+\caption{}+ command.  One can then use the
\verb+\ref{#label}+ command to cross-reference them. As an example, consider
the label declared for Figure \ref{fig1} which is
\verb+\label{fig1}+. To cross-reference it, use the command
\verb+ Figure \ref{fig1}+, for which it comes up as
``Figure~\ref{fig1}".

\subsection{Details on reference citations}\label{subsec3}



\noindent
The following is an example for \verb+\cite{...}+: \cite{haque2025llms}. Another example for \verb+\citep{...}+: \citep{krishnapriya2025hybrid,murphy2012machine,ji20123d}.
Sample cites here \cite{krizhevsky2012imagenet} and \cite{lecun2015deep,ravi2016deep,woods2005open}.


\section{Lists}\label{sec9}

List in \LaTeX{} can be of three types: numbered, bulleted and unnumbered. The ``enumerate'' environment produces a numbered list, the 
``itemize'' environment produces a bulleted list and the ``unlist''
environment produces an unnumbered list.
In each environment, a new entry is added via the \verb+\item+ command.
\begin{enumerate}[1.]
\item This is the 1st item

\item Enumerate creates numbered lists, itemize creates bulleted lists and
unnumerate creates unnumbered lists.
\begin{enumerate}[(a)]
\item Second level numbered list. Enumerate creates numbered lists, itemize creates bulleted lists and
description creates unnumbered lists.

\item Second level numbered list. Enumerate creates numbered lists, itemize creates bulleted lists and
description creates unnumbered lists.
\begin{enumerate}[(ii)]
\item Third level numbered list. Enumerate creates numbered lists, itemize creates bulleted lists and
description creates unnumbered lists.

\item Third level numbered list. Enumerate creates numbered lists, itemize creates bulleted lists and
description creates unnumbered lists.
\end{enumerate}

\item Second level numbered list. Enumerate creates numbered lists, itemize creates bulleted lists and
description creates unnumbered lists.

\item Second level numbered list. Enumerate creates numbered lists, itemize creates bulleted lists and
description creates unnumbered lists.
\end{enumerate}

\item Enumerate creates numbered lists, itemize creates bulleted lists and
description creates unnumbered lists.

\item Numbered lists continue.
\end{enumerate}
Lists in \LaTeX{} can be of three types: enumerate, itemize and description.
In each environment, a new entry is added via the \verb+\item+ command.
\begin{itemize}
\item First level bulleted list. This is the 1st item

\item First level bulleted list. Itemize creates bulleted lists and description creates unnumbered lists.
\begin{itemize}
\item Second level dashed list. Itemize creates bulleted lists and description creates unnumbered lists.

\item Second level dashed list. Itemize creates bulleted lists and description creates unnumbered lists.

\item Second level dashed list. Itemize creates bulleted lists and description creates unnumbered lists.
\end{itemize}

\item First level bulleted list. Itemize creates bulleted lists and description creates unnumbered lists.

\item First level bulleted list. Bullet lists continue.
\end{itemize}

\noindent
Example for unnumbered list items:

\begin{unlist}
\item Sample unnumberd list text. Sample unnumberd list text. Sample unnumberd list text. Sample unnumberd list text. Sample unnumberd list text.

\item Sample unnumberd list text. Sample unnumberd list text. Sample unnumberd list text.

\item sample unnumberd list text. Sample unnumberd list text. Sample unnumberd list text. Sample unnumberd list text. Sample unnumberd list text. Sample unnumberd list text. Sample unnumberd list text.
\end{unlist}

\section{Examples for theorem-like environments}\label{sec10}

For theorem-like environments, we require the \verb+amsthm+ package. There are three types of predefined theorem styles - \verb+thmstyleone+, \verb+thmstyletwo+ and \verb+thmstylethree+   (check your journal's instructions page in case a specific style is required).

\medskip
\noindent\begin{tabular}{|l|p{13pc}|}
\hline
\verb+thmstyleone+ & Numbered, theorem head in bold font and theorem text in italic style \\\hline
\verb+thmstyletwo+ & Numbered, theorem head in roman font and theorem text in italic style \\\hline
\verb+thmstylethree+ & Numbered, theorem head in bold font and theorem text in roman style \\\hline
\end{tabular}


\begin{theorem}[Theorem subhead]\label{thm1}
Example theorem text. Example theorem text. Example theorem text. Example theorem text. Example theorem text.
\end{theorem}

You can write the main content here.You can write the main content here.You can write the main content here.You can write the main content here.

\begin{proposition}
Example proposition text. Example proposition text. Example proposition text. Example proposition text. Example proposition text.
\end{proposition}



\begin{example}
Here is the content of the example.Here is the content of the example.Here is the content of the example.
\end{example}

You can write the main content here.You can write the main content here.You can write the main content here.You can write the main content here.

\begin{remark}
Here is the content of the remark.Here is the content of the remark.Here is the content of the remark.
\end{remark}

You can write the main content here.You can write the main content here.You can write the main content here.You can write the main content here.

\begin{definition}[Definition sub head]
Example definition text. Example definition text. Example definition text. Example definition text. Example definition text. Example definition text. Example definition text. Example definition text.
\end{definition}

Apart from the above styles, we have the \verb+\begin{proof}+ \verb+...+ \verb+\end{proof}+ environment - with the proof head in italic style and the body text in roman font with an open square at the end of each proof environment.

\begin{proof}Example for proof text. Example for proof text. Example for proof text. Example for proof text. Example for proof text. Example for proof text. Example for proof text. Example for proof text. Example for proof text. Example for proof text.
\end{proof}

You can write the main content here.You can write the main content here.You can write the main content here.You can write the main content here.

\begin{proof}[Proof of Theorem~{\upshape\ref{thm1}}]
Example for proof text. Example for proof text. Example for proof text. Example for proof text. Example for proof text. Example for proof text. Example for proof text. Example for proof text. Example for proof text. Example for proof text.
\end{proof}

\noindent
For a quote environment, one has to use\newline \verb+\begin{quote}...\end{quote}+
\begin{quote}
Quoted text example. Quoted text example. Quoted text example. Quoted text example. 
\end{quote}
Here is an example text(refer Figure~\ref{fig3}). Here is an example text(refer Table~\ref{tab3}).

\section{Conclusion}

Some Conclusions here.

%%%%%%%%%%%%%%

\begin{appendices}

\section{Section title of first appendix}\label{sec11}

Here is an example text.Here is an example text.Here is an example text.Here is an example text.Here is an example text.Here is an example text.Here is an example text.

\subsection{Subsection title of first appendix}\label{subsec4}

Here is an example text.Here is an example text.Here is an example text.Here is an example text.Here is an example text.Here is an example text.Here is an example text.

\subsubsection{Subsubsection title of first appendix}\label{subsubsec3}

Example for an unnumbered figure:

\begin{figure}[!h]
\centering
{\color{black!20}\rule{85pt}{92pt}}
\end{figure}

Here is an example text.Here is an example text.Here is an example text.Here is an example text.Here is an example text.Here is an example text.Here is an example text.

\section{Section title of second appendix}\label{sec12}%

Here is an example text.Here is an example text.Here is an example text.Here is an example text.Here is an example text.Here is an example text.Here is an example text.

\begin{figure}[b]
\centering
{\color{black!20}\rule{217pt}{120pt}}
\caption{This is an example for appendix figure\label{fig4}}
\end{figure}

\begin{table}[t]%
\begin{center}
\begin{minipage}{.52\columnwidth}
\caption{This is an example of Appendix table \label{tab4}}%
\begin{tabular}{@{}lcc@{}}%
\toprule
col1 head & col2 head & col3 head \\
\midrule
col1 text & col2 text & col3 text \\
col1 text & col2 text & col3 text \\
col1 text & col2 text & col3 text \\
\botrule
\end{tabular}
\end{minipage}
\end{center}
\end{table}

\subsection{Subsection title of second appendix}\label{subsec5}

Here is an example text.Here is an example text.Here is an example text.Here is an example text.Here is an example text.Here is an example text.Here is an example text.


\subsubsection{Subsubsection title of second appendix}\label{subsubsec4}

Here is an example text.Here is an example text.Here is an example text.Here is an example text.Here is an example text.Here is an example text.Here is an example text.


Example for an equation inside the appendix:
\begin{align}
 p &= \frac{\gamma^{2} - (n_{C} -1)H}{(n_{C} - 1) + H - 2\gamma}, \label{1eq:hybobo:pfromgH} \\
 \theta &= \frac{(\gamma - H)^{2}(\gamma - n_{C} -1)^{2}}{(n_{C} - 1 + H - 2\gamma)^{2}} \label{2eq:hybobo:tfromgH}\; .
\end{align}

\section{Example of another appendix section}\label{sec13}%

Here is an example text.Here is an example text.Here is an example text.Here is an example text.Here is an example text.Here is an example text.Here is an example text.
\begin{equation}
\mathcal{L} = i \bar{\psi} \gamma^\mu D_\mu \psi
    - \frac{1}{4} F_{\mu\nu}^a F^{a\mu\nu} - m \bar{\psi} \psi.
\label{eq26}
\end{equation}

Here is an example text.Here is an example text.Here is an example text.Here is an example text.Here is an example text.Here is an example text.Here is an example text.


Here is an example text.Here is an example text.Here is an example text.Here is an example text.Here is an example text.Here is an example text.Here is an example text.

%% Example for unnumbered table inside appendix
\begin{table}
\begin{center}
\begin{minipage}{.52\columnwidth}
\caption{}{%
\begin{tabular}{lcc}%
\toprule
col1 head & col2 head & col3 head \\
\midrule
col1 text & col2 text & col3 text \\
col1 text & col2 text & col3 text \\
col1 text & col2 text & col3 text \\
\botrule
\end{tabular}}{}
\end{minipage}
\end{center}
\end{table}

\end{appendices}


\section{Acknowledgments}
The authors thank the anonymous reviewers for their valuable suggestions. This work is supported by the national research programs. Experiments utilized institutional computing infrastructure. Methodological refinements benefited from expert reviewer commentary.

\bibliographystyle{TBench}
\bibliography{reference}
\end{document}
